\section{Utilisation des outils de suivi de projet}

Cette section présente les outils mis en place pour assurer la qualité et la traçabilité du code tout au long du développement.

\subsection{GitLab CI/CD : Intégration et déploiement continus}

L'intégration continue (CI) et le déploiement continu (CD) sont au cœur de notre processus de développement. GitLab CI/CD automatise la compilation, les tests et le déploiement du projet.

\subsubsection{Pipeline de CI/CD}

Le pipeline GitLab est organisé en plusieurs stages séquentiels :

\begin{enumerate}
    \item \textbf{build} : Compilation du code source
    \item \textbf{test} : Exécution des tests unitaires et d'intégration
    \item \textbf{sonarqube-check} : Analyse statique du code
    \item \textbf{package} : Création des artefacts JAR
    \item \textbf{deploy} : Déploiement sur Nexus
    \item \textbf{release} : Gestion des versions
\end{enumerate}

\subsubsection{Stratégie de tests automatisés}

Chaque modification de code déclenche automatiquement :
\begin{itemize}
    \item L'exécution de tous les tests unitaires
    \item La génération des rapports de couverture de code (JaCoCo)
    \item La validation de la couverture minimale (80\%)
    \item Les tests d'intégration inter-modules
\end{itemize}

\subsubsection{Qualité de code et métriques}

Le pipeline intègre des vérifications de qualité :
\begin{itemize}
    \item Vérification du style de code
    \item Détection des vulnérabilités
    \item Analyse de la complexité cyclomatique
    \item Rapport de dette technique
\end{itemize}

\subsection{Nexus Repository Manager : Gestion des artefacts}

Nexus est utilisé comme gestionnaire de dépôts Maven pour stocker et distribuer les artefacts du projet.

\subsubsection{Organisation des repositories}

\begin{itemize}
    \item \textbf{snapshots} : Versions de développement (suffixe -SNAPSHOT)
    \item \textbf{releases} : Versions stables et versionnées
    \item \textbf{proxy} : Cache des dépendances Maven Central
\end{itemize}

\subsubsection{Workflow de release}

Le processus de release utilise le plugin Maven Release :

\begin{enumerate}
    \item Préparation : incrémentation de version, création du tag Git
    \item Vérification : tests automatiques sur la version candidate
    \item Publication : déploiement sur le repository releases
    \item Archivage : conservation des artefacts et métadonnées
\end{enumerate}

\subsubsection{Avantages constatés}

L'utilisation de Nexus apporte plusieurs bénéfices :
\begin{itemize}
    \item \textbf{Traçabilité} : Historique complet des versions déployées
    \item \textbf{Performance} : Cache des dépendances externes (réduction du temps de build)
    \item \textbf{Fiabilité} : Isolation des builds par rapport aux pannes de Maven Central
    \item \textbf{Gouvernance} : Contrôle des dépendances utilisées dans le projet
\end{itemize}

\subsection{SonarQube : Analyse statique et qualité du code}

SonarQube réalise une analyse approfondie du code source pour détecter les bugs, vulnérabilités et mauvaises pratiques.

\subsubsection{Métriques suivies}

\begin{itemize}
    \item \textbf{Couverture de code} : Pourcentage de code couvert par les tests
    \item \textbf{Duplication} : Détection du code dupliqué
    \item \textbf{Complexité} : Mesure de la complexité cyclomatique
    \item \textbf{Maintenabilité} : Note globale de maintenabilité (A-E)
    \item \textbf{Fiabilité} : Nombre et sévérité des bugs détectés
    \item \textbf{Sécurité} : Vulnérabilités et points faibles de sécurité
\end{itemize}

\subsubsection{Quality Gates}

Des seuils de qualité sont définis :
\begin{itemize}
    \item Couverture minimale : 80\%
    \item Aucun bug de sévérité critique ou bloquante
    \item Pas de vulnérabilités de sécurité
    \item Taux de duplication < 3\%
    \item Note de maintenabilité minimale : B
\end{itemize}

\subsubsection{Intégration dans le workflow}

SonarQube est intégré au pipeline CI/CD :
\begin{enumerate}
    \item Analyse automatique à chaque merge request
    \item Blocage du merge si les quality gates ne sont pas satisfaits
    \item Rapport de progression visible dans GitLab
    \item Historique de l'évolution de la qualité
\end{enumerate}

\subsection{Outils de planification et communication}

Le suivi de projet s'appuie également sur des outils collaboratifs.

\subsubsection{Jira : Planification Agile}

Jira a été utilisé pour la gestion des Sprints et le suivi des tâches :
\begin{itemize}
    \item Fixation des objectifs de chaque Sprint
    \item Décomposition des \textbf{Epics} en tâches détaillées (tickets)
    \item Suivi de l'avancement via les burndown charts
    \item Estimation de la complexité via les points d'effort
\end{itemize}

\subsubsection{Discord : Communication d'équipe}

Discord a servi de support aux réunions de groupe régulières et au \textit{daily meeting} pour le suivi de l'avancement. Les canaux dédiés ont permis une communication asynchrone structurée entre les membres de l'équipe.

\subsubsection{Merge Requests et Code Review}

Les Merge Requests (MR) avec revue par les pairs ont été systématiquement utilisées pour garantir la qualité du code. Cela justifie l'écart constaté : les tâches terminées le dimanche sont souvent \textit{reviewées} et \textit{mergées} que le lundi, d'où leur statut inachevé le dernier jour du Sprint.

\subsection{Bilan de l'utilisation des outils}

L'ensemble de ces outils forme un écosystème cohérent qui garantit :

\begin{itemize}
    \item \textbf{Automatisation} : Réduction des tâches manuelles répétitives
    \item \textbf{Qualité} : Détection précoce des régressions et anomalies
    \item \textbf{Collaboration} : Visibilité partagée sur l'état du projet
    \item \textbf{Confiance} : Déploiements fiables et reproductibles
\end{itemize}

Cette infrastructure technique constitue un pilier essentiel de la méthodologie Agile du projet, permettant des itérations rapides tout en maintenant un haut niveau de qualité.
